\chapter{Совместимость полного полевого носителя с determinant на базовом пространстве}
\label{ch:v8-full-42-carrier-base-k-determinant-compatibility-gate}

Ранний determinant использовал три скалярные моды и один массивный
калибровочный канал. На нынешнем полном носителе вакуумная грамова часть и
коммутаторы с фоном инцидентности дают вместо этого
\begin{equation}
 \rank H_{\Phi}=28,\quad\dim\ker H_{\Phi}=2,
 \qquad \rank H_B=3,\quad\dim\ker H_B=9.
 \label{eq:v8-full-carrier-vacuum-ranks}
\end{equation}
После нормировки кинетических блоков отношением $3:1$ точные четвёртые
моменты массовых операторов равны
\begin{equation}
 \Tr(M_{\Phi}^4)=\frac{23053}{18},
 \qquad \Tr(M_B^4)=\frac{36897}{722}.
 \label{eq:v8-full-carrier-fourth-moments}
\end{equation}
Босонный числитель коэффициента Коулмана--Вайнберга поэтому равен
\begin{equation}
 N_{\rm bos}=\Tr(M_{\Phi}^4)+3\Tr(M_B^4)
 =\frac{4659176}{3249}.
 \label{eq:v8-full-carrier-bosonic-ledger}
\end{equation}
Он не совпадает с ранним числом $67$. Для конечного фона точно
\begin{equation}
 \Tr D_F^4=46,
 \qquad N_{\rm full}=N_{\rm bos}-\nu_f\Tr D_F^4,
 \label{eq:v8-full-carrier-fermion-boundary}
\end{equation}
но физическая кратность $\nu_f$ и полный BV-фактор ещё не выведены.

\begin{theorem}[Частичный перенос determinant-ledger]
Полный 42-мерный носитель допускает точный босонный base-$K$ ledger. Он
содержит 28 массивных скалярных и три массивных калибровочных направления и
не является продолжением раннего малого ledger простым умножением кратности.
\end{theorem}

\begin{nogocriterion}[Полный коэффициент ещё не определён]
Число $N_{\rm bos}$ нельзя называть полным $B$: остаются две скалярные
нулевые моды, девять безмассовых калибровочных направлений, BV-факторизация
и перевод конечного следа $\Tr D_F^4=46$ в физическую фермионную кратность.
\end{nogocriterion}

\begin{protocol}
\begin{enumerate}
 \item массивных скалярных направлений: \textbf{$28$};
 \item массивных калибровочных направлений: \textbf{$3$};
 \item ненарушенных калибровочных направлений: \textbf{$9$};
 \item босонный четвёртый момент: \textbf{$4659176/3249$};
 \item ранний числитель $67$: \textbf{не переносится};
 \item полный суперслед: \textbf{не закрыт};
 \item следующий гейт: \textbf{полный BV-вакуумный фактор и фермионная кратность}.
\end{enumerate}
\end{protocol}

Машинный сертификат сохранён в
\path{s2t_v8_full_42_carrier_base_k_determinant_compatibility_gate_results.json}.